\documentclass{amsart}
% For dealing with references we use the comment environment
\usepackage{verbatim}
\newenvironment{reference}{\comment}{\endcomment}
%\newenvironment{reference}{}{}
\newenvironment{slogan}{\comment}{\endcomment}
\newenvironment{history}{\comment}{\endcomment}

% For commutative diagrams we use Xy-pic
\usepackage[all]{xy}

% We use 2cell for 2-commutative diagrams.
\xyoption{2cell}
\UseAllTwocells

% We use multicol for the list of chapters between chapters
\usepackage{multicol}

% This is generally recommended for better output
\usepackage{lmodern}
\usepackage[T1]{fontenc}

% For cross-file-references

% Package for hypertext links:
\usepackage{hyperref}

% For any local file, say "hello.tex" you want to link to please

% Theorem environments.
%
\theoremstyle{plain}
\newtheorem{theorem}[subsection]{Theorem}
\newtheorem{proposition}[subsection]{Proposition}
\newtheorem{lemma}[subsection]{Lemma}

\theoremstyle{definition}
\newtheorem{definition}[subsection]{Definition}
\newtheorem{example}[subsection]{Example}
\newtheorem{exercise}[subsection]{Exercise}
\newtheorem{situation}[subsection]{Situation}

\theoremstyle{remark}
\newtheorem{remark}[subsection]{Remark}
\newtheorem{remarks}[subsection]{Remarks}

\numberwithin{equation}{subsection}

% Macros
%
\def\lim{\mathop{\mathrm{lim}}\nolimits}
\def\colim{\mathop{\mathrm{colim}}\nolimits}
\def\Spec{\mathop{\mathrm{Spec}}}
\def\Hom{\mathop{\mathrm{Hom}}\nolimits}
\def\Ext{\mathop{\mathrm{Ext}}\nolimits}
\def\SheafHom{\mathop{\mathcal{H}\!\mathit{om}}\nolimits}
\def\SheafExt{\mathop{\mathcal{E}\!\mathit{xt}}\nolimits}
\def\Sch{\mathit{Sch}}
\def\Mor{\mathop{\mathrm{Mor}}\nolimits}
\def\Ob{\mathop{\mathrm{Ob}}\nolimits}
\def\Sh{\mathop{\mathit{Sh}}\nolimits}
\def\NL{\mathop{N\!L}\nolimits}
\def\CH{\mathop{\mathrm{CH}}\nolimits}
\def\proetale{{pro\text{-}\acute{e}tale}}
\def\etale{{\acute{e}tale}}
\def\QCoh{\mathit{QCoh}}
\def\Ker{\mathop{\mathrm{Ker}}}
\def\Im{\mathop{\mathrm{Im}}}
\def\Coker{\mathop{\mathrm{Coker}}}
\def\Coim{\mathop{\mathrm{Coim}}}

% Boxtimes
%
\DeclareMathSymbol{\boxtimes}{\mathbin}{AMSa}{"02}

%
% Macros for moduli stacks/spaces
%
\def\QCohstack{\mathcal{QC}\!\mathit{oh}}
\def\Cohstack{\mathcal{C}\!\mathit{oh}}
\def\Spacesstack{\mathcal{S}\!\mathit{paces}}
\def\Quotfunctor{\mathrm{Quot}}
\def\Hilbfunctor{\mathrm{Hilb}}
\def\Curvesstack{\mathcal{C}\!\mathit{urves}}
\def\Polarizedstack{\mathcal{P}\!\mathit{olarized}}
\def\Complexesstack{\mathcal{C}\!\mathit{omplexes}}
% \Pic is the operator that assigns to X its picard group, usage \Pic(X)
% \Picardstack_{X/B} denotes the Picard stack of X over B
% \Picardfunctor_{X/B} denotes the Picard functor of X over B
\def\Pic{\mathop{\mathrm{Pic}}\nolimits}
\def\Picardstack{\mathcal{P}\!\mathit{ic}}
\def\Picardfunctor{\mathrm{Pic}}
\def\Deformationcategory{\mathcal{D}\!\mathit{ef}}

\setlength{\emergencystretch}{3em}
\begin{document}
\title[Stacks Project, Tag 06IW]{399 --- Existence of versal formal objects under Schlessinger conditions}
\maketitle
\noindent Copyright (C) 2005--2025 Johan de Jong. Stacks Project authors. Source: \href{https://stacks.math.columbia.edu/tag/06IW}{Tag 06IW}.

\noindent Editorial difficulty note (not author text). Difficulty H2: The proof constructs a descending sequence of minimal lifting ideals inside a power-series ring and takes its complete quotient. It then proves versality by analyzing an arbitrary small-extension pullback, distinguishing split and essential surjections, and forcing its kernel to contain the next minimal ideal. The inductive construction and the universal lifting argument are both substantive..

\noindent Editorial provenance note (not author text). History: excerpt from revision \texttt{a0444\allowbreak 6e57e\allowbreak c1fbc\allowbreak 252a8\allowbreak 71afc\allowbreak ec775\allowbreak 2fb28\allowbreak 07b14}, formal-defos.tex, lines 3866--4040. Reference presentation was adapted; original-author context and core support blocks were retained or added. The mathematical wording is unchanged. Licensed under GNU FDL 1.2 or later; no invariant sections or cover texts. See ../licenses/COPYING. Context blocks reproduce author definitions and supporting results; they are not additional counted problems.

\begin{definition}
\label{definition-CLambda}
Let $\Lambda$ be a Noetherian ring and let $\Lambda \to k$ be a finite
ring map where $k$ is a field. We define {\it $\mathcal{C}_\Lambda$} to be
the category with
\begin{enumerate}
\item objects are pairs $(A, \varphi)$ where $A$ is an Artinian local
$\Lambda$-algebra and where $\varphi : A/\mathfrak m_A \to k$ is a
$\Lambda$-algebra isomorphism, and
\item morphisms $f : (B, \psi) \to (A, \varphi)$ are local $\Lambda$-algebra
homomorphisms such that $\varphi \circ (f \bmod \mathfrak m) = \psi$.
\end{enumerate}
We say we are in the {\it classical case} if $\Lambda$ is a Noetherian
complete local ring and $k$ is its residue field.
\end{definition}

\section{Notation and Conventions}
\label{section-notations-conventions}

\noindent
A ring is commutative with $1$. The  maximal ideal of a local ring $A$
is denoted by $\mathfrak{m}_A$. The set of positive integers is denoted
by $\mathbf{N} = \{1, 2, 3, \ldots\}$. If $U$ is an object of a
category $\mathcal{C}$, we denote by $\underline{U}$
the functor
$\Mor_\mathcal{C}(U, -): \mathcal{C} \to \textit{Sets}$, see
Remarks \href{https://stacks.math.columbia.edu/tag/06GK}{remarks cofibered groupoids (Tag 06GK)} (\href{https://stacks.math.columbia.edu/tag/06GQ}{item definition yoneda (Tag 06GQ)}).
Warning: this may conflict with the notation in other chapters where we
sometimes use $\underline{U}$ to denote $h_U(-) = \Mor_\mathcal{C}(-, U)$.

\medskip\noindent
Throughout this chapter $\Lambda$ is a Noetherian ring and
$\Lambda \to k$ is a finite ring map from $\Lambda$ to a field.
The kernel of this map is denoted $\mathfrak m_\Lambda$ and the
image $k' \subset k$. It turns out that $\mathfrak m_\Lambda$ is
a maximal ideal, $k' = \Lambda/\mathfrak m_\Lambda$ is a field, and
the extension $k/k'$ is finite. See discussion surrounding
(\href{https://stacks.math.columbia.edu/tag/06S2}{equation k prime (Tag 06S2)}).




\begin{definition}
\label{definition-completion-CLambda}
Let $\Lambda$ be a Noetherian ring and let $\Lambda \to k$ be a finite
ring map where $k$ is a field. We define {\it $\widehat{\mathcal{C}}_\Lambda$}
to be the category with
\begin{enumerate}
\item objects are pairs $(R, \varphi)$ where $R$ is a Noetherian complete
local $\Lambda$-algebra and where $\varphi : R/\mathfrak m_R \to k$ is a
$\Lambda$-algebra isomorphism, and
\item morphisms $f : (S, \psi) \to (R, \varphi)$ are local $\Lambda$-algebra
homomorphisms such that $\varphi \circ (f \bmod \mathfrak m) = \psi$.
\end{enumerate}
\end{definition}

\begin{definition}
\label{definition-category-cofibred-groupoids}
Let $\mathcal{C}$ be a category.  A {\it category cofibered in groupoids over
$\mathcal{C}$} is a category $\mathcal{F}$ equipped with a functor
$p: \mathcal{F} \to \mathcal{C}$ such that $\mathcal{F}^{opp}$ is a category
fibered in groupoids over $\mathcal{C}^{opp}$ via
$p^{opp}: \mathcal{F}^{opp} \to \mathcal{C}^{opp}$.
\end{definition}

\begin{definition}
\label{definition-predeformation-category}
A {\it predeformation category} $\mathcal{F}$ is a category cofibered
in groupoids over $\mathcal{C}_\Lambda$ such that $\mathcal{F}(k)$ is
equivalent to a category with a single object and a single morphism,
i.e., $\mathcal{F}(k)$ contains at least one object and there is a
unique morphism between any two objects. A {\it morphism of predeformation
categories} is a morphism of categories cofibered in groupoids over
$\mathcal{C}_\Lambda$.
\end{definition}

\begin{definition}
\label{definition-formal-objects}
Let $\mathcal{F}$ be a category cofibered in groupoids over
$\mathcal{C}_\Lambda$. The {\it category $\widehat{\mathcal{F}}$ of formal
objects of  $\mathcal{F}$} is the category with the following objects and
morphisms.
\begin{enumerate}
\item A {\it formal object $\xi = (R, \xi_n, f_n)$ of $\mathcal{F}$}
consists of an object $R$ of $\widehat{\mathcal{C}}_\Lambda$, and a collection
indexed by $n \in \mathbf{N}$ of objects $\xi_n$ of
$\mathcal{F}(R/\mathfrak m_R^n)$ and morphisms
$f_n : \xi_{n + 1} \to \xi_n$ lying over the projection
$R/\mathfrak m_R^{n + 1} \to R/\mathfrak m_R^n$.
\item Let $\xi = (R, \xi_n, f_n)$ and $\eta = (S, \eta_n, g_n)$ be
formal objects of $\mathcal{F}$.  A {\it morphism $a : \xi \to \eta$ of
formal objects} consists of a map $a_0 : R \to S$ in
$\widehat{\mathcal{C}}_\Lambda$ and a collection $a_n : \xi_n \to \eta_n$
of morphisms of $\mathcal{F}$ lying over
$R/\mathfrak m_R^n \to S/\mathfrak m_S^n$,
such that for every $n$ the diagram
$$
\xymatrix{
\xi_{n + 1} \ar[r]_{f_n} \ar[d]_{a_{n + 1}} & \xi_n \ar[d]^{a_n} \\
\eta_{n + 1} \ar[r]^{g_n} & \eta_n
}
$$
commutes.
\end{enumerate}
\end{definition}

\begin{definition}
\label{definition-versal}
Let $\mathcal{F}$ be a category cofibered in groupoids.  Let $\xi$ be a formal
object of $\mathcal{F}$ lying over $R \in \Ob(\widehat{\mathcal{C}}_\Lambda)$.
We say $\xi$ is {\it versal} if the corresponding morphism
$\underline{\xi}: \underline{R}|_{\mathcal{C}_\Lambda} \to \mathcal{F}$
of Remark \href{https://stacks.math.columbia.edu/tag/06HC}{remark formal objects yoneda (Tag 06HC)} is smooth.
\end{definition}

\begin{definition}
\label{definition-S1-S2}
Let $\mathcal{F}$ be a category cofibered in groupoids over $\mathcal
C_\Lambda$. We define {\it conditions (S1) and (S2)}
on $\mathcal{F}$ as follows:
\begin{enumerate}
\item[(S1)] Every diagram in $\mathcal{F}$
$$
\vcenter{
\xymatrix{
           & x_2 \ar[d] \\
x_1 \ar[r] & x
}
}
\quad\text{lying over}\quad
\vcenter{
\xymatrix{
           & A_2 \ar[d] \\
A_1 \ar[r] & A
}
}
$$
in $\mathcal{C}_\Lambda$ with $A_2 \to A$ surjective can be completed
to a commutative diagram
$$
\vcenter{
\xymatrix{
y \ar[r] \ar[d] & x_2 \ar[d] \\
x_1 \ar[r]      & x
}
}
\quad\text{lying over}\quad
\vcenter{
\xymatrix{
A_1 \times_A A_2 \ar[r] \ar[d] & A_2 \ar[d] \\
A_1 \ar[r]      & A.
}
}
$$
\item[(S2)]
The condition of (S1) holds for diagrams in $\mathcal{F}$ lying over
a diagram in $\mathcal{C}_\Lambda$ of the form
$$
\xymatrix{
          & k[\epsilon] \ar[d] \\
A  \ar[r] & k.
}
$$
Moreover, if we have two commutative diagrams in $\mathcal{F}$
$$
\vcenter{
\xymatrix{
y \ar[r]_c \ar[d]_a & x_\epsilon \ar[d]^e \\
x \ar[r]^d          & x_0
}
}
\quad\text{and}\quad
\vcenter{
\xymatrix{
y' \ar[r]_{c'} \ar[d]_{a'} & x_\epsilon \ar[d]^e \\
x \ar[r]^d                 & x_0
}
}
\quad\text{lying over}\quad
\vcenter{
\xymatrix{
A \times_k k[\epsilon] \ar[r] \ar[d] & k[\epsilon] \ar[d] \\
A  \ar[r] & k
}
}
$$
then there exists a morphism $b : y \to y'$ in
$\mathcal{F}(A \times_k k[\epsilon])$ such that $a = a' \circ b$.
\end{enumerate}
\end{definition}

\begin{definition}
\label{definition-tangent-space}
Let $\mathcal{F}$ be a predeformation category.
The {\it tangent space $T \mathcal{F}$ of $\mathcal{F}$}
is the set $\overline{\mathcal{F}}(k[\epsilon])$
of isomorphism classes of objects in the fiber category $\mathcal
F(k[\epsilon])$.
\end{definition}

\begin{definition}
\label{definition-small-extension}
Let $f: B \to A$ be a ring map in $\mathcal{C}_\Lambda$.  We say $f$
is a {\it small extension} if it is surjective and $\Ker(f)$ is a nonzero
principal ideal which is annihilated by $\mathfrak{m}_B$.
\end{definition}

\noindent
Explicitly, $p: \mathcal{F} \to \mathcal{C}$ is cofibered in groupoids if
the following two conditions hold:
\begin{enumerate}
\item For every morphism $f: U \to V$ in $\mathcal{C}$ and every object
$x$ lying over $U$, there is a morphism $x \to y$ of $\mathcal{F}$ lying
over $f$.
\item For every pair of morphisms $a: x \to y$ and $b: x \to z$
of $\mathcal{F}$ and any morphism $f: p(y) \to p(z)$ such that $p(b) = f
\circ p(a)$, there exists a unique morphism $c: y \to z$ of $\mathcal
F$ lying over $f$ such that $b = c \circ a$.
\end{enumerate}


\medskip\noindent
We denote by $k[\epsilon]$ the ring of dual numbers over $k$.  More
generally, for a $k$-vector space $V$, we denote by $k[V]$ the $k$-algebra
whose underlying vector space is $k \oplus V$ and whose multiplication is given
by $(a, v) \cdot (a', v') = (aa', av' + a'v)$.  When $V = k$, $k[V]$ is the ring
of dual numbers over $k$.  For any finite dimensional $k$-vector space $V$
the ring $k[V]$ is in $\mathcal{C}_\Lambda$.


\begin{lemma}
\label{lemma-versal-object-existence}
Let $\mathcal{F}$ be a category cofibred in groupoids over
$\mathcal{C}_\Lambda$. Assume the following conditions hold:
\begin{enumerate}
\item $\mathcal{F}$ is a predeformation category.
\item $\mathcal{F}$ satisfies (S1).
\item $\mathcal{F}$ satisfies (S2).
\item $\dim_k T\mathcal{F}$ is finite.
\end{enumerate}
Then $\mathcal{F}$ has a versal formal object.
\end{lemma}

\begin{proof}
Assume (1), (2), (3), and (4) hold. Choose an object
$R \in \Ob(\widehat{\mathcal{C}}_\Lambda)$
such that $\underline{R}|_{\mathcal{C}_\Lambda}$ is smooth.
See Lemma \href{https://stacks.math.columbia.edu/tag/06SM}{lemma exists smooth (Tag 06SM)}.
Let $r = \dim_k T\mathcal{F}$ and put $S = R[[X_1, \ldots, X_r]]$.

\medskip\noindent
We are going to inductively construct for $n \geq 2$ pairs
$(J_n, f_{n - 1} : \xi_n \to \xi_{n - 1})$
where $J_n \subset S$ is an decreasing sequence of ideals and
$f_{n - 1} : \xi_n \to \xi_{n - 1}$ is a morphism of
$\mathcal{F}$ lying over the projection $S/J_n \to S/J_{n - 1}$.

\medskip\noindent
Step 1. Let $J_1 = \mathfrak m_S$. Let $\xi_1$ be the unique
(up to unique isomorphism) object of $\mathcal{F}$ over
$k = S/J_1 = S/\mathfrak m_S$

\medskip\noindent
Step 2. Let
$J_2 = \mathfrak m_S^2 + \mathfrak{m}_R S$. Then
$S/J_2 = k[V]$ with $V = kX_1 \oplus \ldots \oplus kX_r$
By (S2) for $\overline{\mathcal{F}}$ we get a bijection
$$
\overline{\mathcal{F}}(S/J_2)
\longrightarrow
T\mathcal{F} \otimes_k V,
$$
see
Lemmas \href{https://stacks.math.columbia.edu/tag/06I0}{lemma S1 S2 associated functor (Tag 06I0)} and
\href{https://stacks.math.columbia.edu/tag/06IH}{lemma tangent space vector space (Tag 06IH)}.
Choose a basis $\theta_1, \ldots, \theta_r$ for $T\mathcal{F}$ and set
$\xi_2 = \sum \theta_i \otimes X_i \in \Ob(\mathcal{F}(S/J_2))$.
The point of this choice is that
$$
d\xi_2 :
\Mor_{\mathcal{C}_\Lambda}(S/J_2, k[\epsilon])
\longrightarrow
T\mathcal{F}
$$
is surjective. Let $f_1 : \xi_2 \to \xi_1$ be the unique morphism.

\medskip\noindent
Induction step. Assume $(J_n, f_{n - 1} : \xi_n \to \xi_{n - 1})$ has been
constructed for some $n \geq 2$. There is a minimal element $J_{n + 1}$ of
the set of ideals $J \subset S$ satisfying:
(a) $\mathfrak m_S J_n \subset J \subset J_n$ and
(b) there exists a morphism $\xi_{n + 1} \to \xi_n$ lying over
$S/J \to S/J_n$, see
Lemma \href{https://stacks.math.columbia.edu/tag/06SZ}{lemma largest closed where lift (Tag 06SZ)}.
Let $f_n : \xi_{n + 1} \to \xi_n$ be any morphism of $\mathcal{F}$
lying over $S/J_{n + 1} \to S/J_n$.

\medskip\noindent
Set $J = \bigcap J_n$. Set $\overline{S} = S/J$. Set $\overline{J}_n = J_n/J$.
By
Lemma \href{https://stacks.math.columbia.edu/tag/06SE}{lemma m adic topology (Tag 06SE)}
the sequence of ideals $(\overline{J}_n)$ induces the
$\mathfrak m_{\overline{S}}$-adic topology on $\overline{S}$.
Since $(\xi_n, f_n)$ is an object of
$\widehat{\mathcal{F}}_\mathcal{I}(\overline{S})$, where $\mathcal{I}$
is the filtration $(\overline{J}_n)$ of $\overline{S}$,
we see that $(\xi_n, f_n)$
induces an object $\xi$ of $\widehat{\mathcal{F}}(\overline{S})$.
see
Lemma \href{https://stacks.math.columbia.edu/tag/06H6}{lemma formal objects different filtration (Tag 06H6)}.

\medskip\noindent
We prove $\xi$ is versal. For versality it suffices to check
conditions (1) and (2) of
Lemma \href{https://stacks.math.columbia.edu/tag/06IU}{lemma versal criterion (Tag 06IU)}.
Condition (1) follows from our choice of $\xi_2$ in Step 2 above.
Suppose given a diagram in $\widehat{\mathcal{F}}$
$$
\vcenter{
\xymatrix{
            &  y \ar[d] \\
\xi \ar[r]  &  x
}
}
\quad\text{lying over}\quad
\vcenter{
\xymatrix{
         &   B  \ar[d]^{f} \\
\overline{S} \ar[r] &   A
}
}
$$
in $\widehat{\mathcal{C}}_\Lambda$ with $f: B \to A$ a small extension
of Artinian rings. We have to show there is a map $\overline{S} \to B$ fitting
into the diagram on the right. Choose $n$ such that
$\overline{S} \to A$ factors through $\overline{S} \to S/J_n$. This is
possible as the sequence $(\overline{J}_n)$ induces the
$\mathfrak m_{\overline{S}}$-adic topology as we saw above.
The pushforward of $\xi$ along $\overline{S} \to S/J_n$ is $\xi_n$.
We may factor $\xi \to x$ as $\xi \to \xi_n \to x$ hence we get a diagram
in $\mathcal{F}$
$$
\vcenter{
\xymatrix{
            &  y \ar[d] \\
\xi_n \ar[r]  &  x
}
}
\quad\text{lying over}\quad
\vcenter{
\xymatrix{
         &   B  \ar[d]^{f} \\
S/J_n \ar[r] &   A .
}
}
$$
To check condition (2) of
Lemma \href{https://stacks.math.columbia.edu/tag/06IU}{lemma versal criterion (Tag 06IU)}
it suffices to complete the diagram
$$
\xymatrix{
S/J_{n + 1} \ar[d] \ar@{-->}[r] & B \ar[d]^{f} \\
S/J_n   \ar[r] & A
}
$$
or equivalently, to complete the diagram
$$
\xymatrix{
  &  S/J_n \times_A B \ar[d]^{p_1} \\
S/J_{n + 1} \ar@{-->}[ur] \ar[r] &  S/J_n.
}
$$
If $p_1$ has a section we are done. If not, by
Lemma \href{https://stacks.math.columbia.edu/tag/06GH}{lemma fiber product CLambda (Tag 06GH)} (2)
$p_1$ is a small extension, so by
Lemma \href{https://stacks.math.columbia.edu/tag/06H0}{lemma essential surjection (Tag 06H0)} (4)
$p_1$ is an essential surjection. Recall that $S = R[[X_1, \ldots, X_r]]$
and that we chose $R$ such that $\underline{R}|_{\mathcal{C}_\Lambda}$
is smooth. Hence there exists a map $h : R \to B$ lifting the map
$R \to S \to S/J_n \to A$. By the universal property of a power series
ring there is an $R$-algebra map $h : S = R[[X_1, \ldots, X_r]] \to B$
lifting the given map $S \to S/J_n \to A$. This induces a map
$g: S \to S/J_n \times_A B$ making the solid square in the diagram
$$
\xymatrix{
S \ar[d] \ar[r]_-g  &  S/J_n \times_A B \ar[d]^{p_1} \\
S/J_{n + 1} \ar@{-->}[ur] \ar[r] &  S/J_n
}
$$
commute. Then $g$ is a surjection since $p_1$ is an essential surjection.
We claim the ideal $K = \Ker(g)$ of $S$ satisfies conditions (a) and
(b) of the construction of $J_{n + 1}$ in the induction step above.
Namely, $K \subset J_n$ is clear and $\mathfrak m_SJ_n \subset K$ as $p_1$
is a small extension; this proves (a). By (S1) applied to
$$
\xymatrix{
            &  y \ar[d] \\
\xi_n \ar[r]  &  x,
}
$$
there exists a lifting of $\xi_n$ to $S/K \cong S/J_n \times_A B$, so (b)
holds. Since $J_{n + 1}$ was the minimal ideal with properties (a) and (b)
this implies $J_{n + 1} \subset K$. Thus the desired map
$S/J_{n+1} \to S/K \cong S/J_n \times_A B$ exists.
\end{proof}
\end{document}
